% Q5: Parallel RLC Natural Response — Four SPST switches (Nilsson P8.11 style)
%
% Topology (left to right on top rail):
%   Is || Ra  [SW1]  L  [SW2]  R(v_o)  [SW3]  C  [SW4]  Rb--Vdc
%
% t < 0: SW1,SW4 closed; SW2,SW3 open
% t = 0: all switches change state → source-free parallel RLC
\documentclass[border=10pt]{standalone}
\usepackage[american voltages, american currents]{circuitikz}
\usepackage{amsmath}
\renewcommand{\familydefault}{\sfdefault}

\begin{document}
\begin{circuitikz}[
    line width=0.8pt,
    every node/.style={font=\sffamily},
    switch label/.style={font=\sffamily\small, above, yshift=4pt},
    switch name/.style={font=\sffamily\small\itshape, below, yshift=-4pt}
]

% === Left section: Is and Ra in parallel ===

% Current source (arrow up)
\draw (0,0) to[I, l=$I_s$] (0,4);

% Top rail to Ra junction
\draw (0,4) -- (2,4);
\fill (2,4) circle (2pt);

% Ra (vertical, parallel to Is)
\draw (2,4) to[R, l_=$R_a$] (2,0);

% Bottom rail under Is and Ra
\draw (0,0) -- (2,0);

% === SW1: closed, opens at t=0 ===
\draw (2,4) -- (3,4)
      to[closing switch, l={\shortstack{$t{=}0$\\(opens)}}, switch name/.append style={}] (5,4);
\node[switch name] at (4,4) {\textit{SW1}};

% Junction for L branch
\draw (5,4) -- (5.5,4);
\fill (5.5,4) circle (2pt);

% L branch (vertical)
\draw (5.5,4) to[L, l_=$L$] (5.5,0);

% Bottom rail extension
\draw (2,0) -- (5.5,0);

% === SW2: open, closes at t=0 ===
\draw (5.5,4) -- (6.5,4)
      to[opening switch, l={\shortstack{$t{=}0$\\(closes)}}] (8.5,4);
\node[switch name] at (7.5,4) {\textit{SW2}};

% Junction for R branch
\draw (8.5,4) -- (9,4);
\fill (9,4) circle (2pt);

% R branch (vertical) with v_o(t) polarity
\draw (9,4) to[R, l_=$R$, v^=$v_o(t)$] (9,0);

% Bottom rail extension
\draw (5.5,0) -- (9,0);

% === SW3: open, closes at t=0 ===
\draw (9,4) -- (10,4)
      to[opening switch, l={\shortstack{$t{=}0$\\(closes)}}] (12,4);
\node[switch name] at (11,4) {\textit{SW3}};

% Junction for C branch
\draw (12,4) -- (12.5,4);
\fill (12.5,4) circle (2pt);

% C branch (vertical)
\draw (12.5,4) to[C, l_=$C$] (12.5,0);

% Bottom rail extension
\draw (9,0) -- (12.5,0);

% === SW4: closed, opens at t=0 ===
\draw (12.5,4) -- (13.5,4)
      to[closing switch, l={\shortstack{$t{=}0$\\(opens)}}] (15.5,4);
\node[switch name] at (14.5,4) {\textit{SW4}};

% === Right section: Rb and Vdc ===
% Rb (horizontal on top rail)
\draw (15.5,4) to[R, l=$R_b$] (18,4);

% Vdc (vertical, + on top)
\draw (18,4) to[V, v_=$V_{dc}$] (18,0);

% Bottom rail extension
\draw (12.5,0) -- (18,0);

\end{circuitikz}
\end{document}
